\textbf{Reading the results.}
\begin{itemize}
\item \textbf{The search space is flat at the top.} Five of six families reach a CV F1 within 0.005 of each other (Table~\ref{tab:automlfam}). The single decision tree is far behind (0.89): one tree cannot exploit 100k sparse, weakly informative features. Three configurations tie at 0.9827, and the default logistic regression of the search (0.9779) is only half a point of F1 below, on a 20k subsample.
\item \textbf{In-distribution, AutoML trades a little recall for fewer false positives.} On the hold-out test, the selected ridge model has a lower FPR (0.65\% vs 0.93\%) and therefore a better precision at 5\% prevalence (0.888 vs 0.848), for a slightly lower F1 (0.989 vs 0.991).
\item \textbf{Out-of-distribution, AutoML is worse on three of the four binary sources.} F1 drops on CEAS\_08 ($-0.015$), Enron ($-0.030$), Ling ($-0.094$) and SpamAssassin ($-0.055$). FPR improves on CEAS\_08 and Enron but worsens on Ling and SpamAssassin. Only the Nazario recall improves (0.619 vs 0.585). The selected model is \emph{more complex} (word bigrams + character 3--5-grams, 150k features, weak regularisation $\alpha=0.044$, training F1 = 1.000): it fits the training sources better, including their artefacts, and transfers worse.
\item \textbf{Conclusion for v0.1.} Optimising an in-distribution objective selects for memorisation of source-specific patterns. This is direct empirical evidence that the selection criterion, not the model family, is the bottleneck. It makes nested LOSO selection (Section~\ref{sec:plan}) the first priority.
\end{itemize}
